\documentclass[preview]{standalone}

\usepackage{amsmath}
\usepackage{amssymb}
\usepackage{bettelini}
\usepackage{stellar}
\usepackage{definitions}
\usepackage{wrapfig}

\begin{document}

\id{measuretheory-riemann-comparison}
\genpage

\section{Comparison with the Riemann Integral}

Let \([a,b]\subseteq\realnumbers\) be equipped with Lebesgue measure. We
compare Riemann integrable functions with
\(\mathcal L^1([a,b])\).

\begin{snippettheorem}{riemann-integral-embeds-into-lebesgue-integral}{Embedding of the Riemann integral into the Lebesgue integral}
    If \(f\) is Riemann integrable on a compact interval, then it is Lebesgue
    integrable there and the two integrals coincide.
\end{snippettheorem}

One tempting but incorrect measurability argument is to approximate a
nonnegative Riemann integrable function by the step functions associated with
lower sums. Even along refinements these functions need not converge
pointwise to \(f\). For example, if
\(f=\chi_{\{\widetilde x\}}\), every lower step function is zero, including
at \(\widetilde x\), although \(f(\widetilde x)=1\). The exceptional set may
even have the cardinality of the continuum while still having measure zero.

\begin{snippettheorem}{lebesgue-criterion-riemann-integrability}{Lebesgue's criterion for Riemann integrability}
    A function \(f\colon[a,b]\fromto\realnumbers\) is Riemann integrable if
    and only if it is bounded and its set \(D_f\) of discontinuity points has
    Lebesgue measure zero.
\end{snippettheorem}

Thus Riemann integrability permits discontinuities on a set of continuum
cardinality. Not every set can occur as the discontinuity set of a function:
such sets are precisely the \(F_\sigma\) sets. In particular, no real
function is discontinuous exactly at every irrational and continuous exactly
at every rational.

\begin{snippetexercise}{function-discontinuous-exactly-on-rationals}{Discontinuities at all rational points}
    Construct a function whose discontinuity set is exactly
    \(\rationalnumbers\). The Dirichlet function is not suitable because it is
    discontinuous everywhere.
\end{snippetexercise}

One solution is Thomae's function: write a rational in lowest terms as
\(p/q\), set \(f(p/q)=1/q\), and set \(f(x)=0\) for irrational \(x\).

\begin{snippetexercise}{discontinuity-sets-are-f-sigma}{Characterization of discontinuity sets}
    Prove that, for every \(f\colon\realnumbers\fromto\realnumbers\), the set
    \(D_f\) is \(F_\sigma\), and conversely that every \(F_\sigma\) set is the
    discontinuity set of some real-valued function.
\end{snippetexercise}

\section{Almost-Everywhere Continuity and Measurability}

\begin{snippetdefinition}{borel-measure-definition}{Borel measure}
    A measure on a topological space is a \emph{Borel measure} if its
    \(\sigma\)-algebra contains the topology.
\end{snippetdefinition}

Its domain need not be exactly the Borel \(\sigma\)-algebra. Let
\((X,\tau)\) be a topological space and
\((X,\mathcal M,\mu)\) a measure space whose measure is Borel. Without
completeness, being continuous outside a null set does not imply
measurability.

\begin{snippetexample}{almost-everywhere-continuity-needs-completeness}{Completeness is necessary}
    Equip \(\realnumbers\) with its Euclidean topology, Borel
    \(\sigma\)-algebra, and the restriction of Lebesgue measure to Borel
    sets. Let \(E\) be the Cantor set and choose a non-Borel subset
    \(F\subseteq E\). Then
    \[
        f=\chi_F
    \]
    is continuous on the open set \(E^c\), whose complement has measure zero,
    but it is not Borel measurable because
    \(f^{-1}((1/2,+\infty))=F\). The obstruction is precisely that the Borel
    restriction of Lebesgue measure is not complete.
\end{snippetexample}

\begin{snippettheorem}{almost-everywhere-continuous-function-measurable}{Almost-everywhere continuity implies measurability for complete Borel measures}
    Let \((X,\tau)\) also be a complete measure space
    \((X,\mathcal M,\mu)\), with \(\mu\) a Borel measure. If
    \(f\colon X\fromto\mathbb F\) is continuous outside a measurable null
    set, then \(f\) is measurable. In particular, this holds if
    \(D_f\in\mathcal M\) and \(\mu(D_f)=0\).
\end{snippettheorem}

Continuity almost everywhere is not the same as being almost everywhere
equal to a continuous function. A jump discontinuity cannot be removed by
changing only the value at the jump. Conversely, the Dirichlet function is
almost everywhere equal to zero but is continuous nowhere.

\begin{snippetproof}{almost-everywhere-continuous-function-measurable-proof}{almost-everywhere-continuous-function-measurable}{Almost-everywhere continuity implies measurability for complete Borel measures}
    Let \(N\in\mathcal M\) be null and suppose that
    \(\restr{f}{X\setminus N}\) is continuous for the subspace topology.
    It is therefore measurable for the induced \(\sigma\)-algebra. For every
    open \(U\subseteq\mathbb F\), there is \(A\in\mathcal M\) such that
    \[
        f^{-1}(U)\cap(X\setminus N)=A\cap(X\setminus N).
    \]
    The remaining part \(f^{-1}(U)\cap N\) is measurable by completeness, so
    \(f^{-1}(U)\in\mathcal M\).
\end{snippetproof}

Now let \(f\) be Riemann integrable on \([a,b]\). Lebesgue's criterion makes
\(f\) measurable. Since it is bounded and the interval has finite measure,
\(f\) is dominated by an integrable constant and hence belongs to
\(\mathcal L^1([a,b])\).

\begin{snippetproof}{riemann-and-lebesgue-integrals-coincide-proof}{riemann-integral-embeds-into-lebesgue-integral}{Equality of the Riemann and Lebesgue integrals}
    Suppose first that \(f\geq0\). A lower Riemann sum is the integral of a
    measurable simple function below \(f\); hence it is at most
    \(\int_{[a,b]}f\,d\mu\). An upper Riemann sum is the integral of a
    measurable simple function above \(f\), so by monotonicity it is at least
    \(\int_{[a,b]}f\,d\mu\). Since Riemann integrability gives, for every
    \(\varepsilon>0\), lower and upper sums whose difference is less than
    \(\varepsilon\), the Lebesgue integral is squeezed between quantities
    converging to the Riemann integral. Thus
    \[
        \int_a^b f(x)\,dx=\int_{[a,b]}f\,d\mu.
    \]
    Applying this result to \(f^+\) and \(f^-\) proves the general case.
\end{snippetproof}

For this reason, the Riemann notation is also used for Lebesgue integrals over
intervals.

\section{Improper Riemann Integrals}

Absolute improper Riemann integrability is sufficient for Lebesgue
integrability.

\begin{snippetproposition}{absolutely-convergent-improper-riemann-integral-lebesgue}{Absolute improper Riemann integrability implies Lebesgue integrability}
    Let \(-\infty\leq a<b\leq+\infty\). If \(|f|\) is improperly Riemann
    integrable on \((a,b)\), then \(f\in\mathcal L^1((a,b))\) and its
    Lebesgue and improper Riemann integrals coincide.
\end{snippetproposition}

\begin{snippetproof}{absolutely-convergent-improper-riemann-integral-lebesgue-proof}{absolutely-convergent-improper-riemann-integral-lebesgue}{Absolute improper Riemann integrability implies Lebesgue integrability}
    \begin{wrapfigure}{r}{6cm}
        \begin{center}
            \begin{tikzpicture}[scale=0.8]
                \draw[->] (0,0) -- (7,0);
                \draw[->] (0,0) -- (0,5);

                \draw[dashed] (1,0) -- (1,5);
                \draw[dashed] (0,1) -- (7,1);

                \draw[thick,domain=1.3:6.4,samples=200,smooth]
                plot (\x,{1 + 1.8/((\x-1)^0.9)});

                \fill (1,0) circle (0pt) node[below] {$a$};
                \fill (6.2,0) circle (0pt) node[below] {$b$};
            \end{tikzpicture}
        \end{center}
    \end{wrapfigure}
    Consider first the endpoint configuration shown in the figure and a
    nonnegative function. Choose a monotone sequence \(x_n\downarrow a\).
    By the proper-integral comparison,
    \[
        \int_{x_n}^bf(x)\,dx
        =\int_{(a,b)}f(x)\chi_{(x_n,b)}(x)\,d\mu.
    \]
    The integrands increase pointwise to \(f\). Hence the monotone convergence
    theorem identifies the improper limit with the Lebesgue integral.
    Applying this to \(|f|\) proves that \(f\in\mathcal L^1\), and applying it
    to \(f^+\) and \(f^-\) proves equality of the integrals. The other finite
    or infinite endpoint configurations follow by choosing increasing compact
    subintervals of \((a,b)\).
    \wrapfill
\end{snippetproof}

\begin{snippetexercise}{parameter-integral-log-over-quadratic}{A parameter-dependent improper integral}
    Consider
    \[
        F(y)=\int_0^{+\infty}\frac{\ln x}{x^2+y}\,dx.
    \]
    Determine its domain and study continuity.
\end{snippetexercise}

\begin{snippetsolution}{parameter-integral-log-over-quadratic-solution}{parameter-integral-log-over-quadratic}{A parameter-dependent improper integral}
    If \(y>0\), the denominator never vanishes. Near \(+\infty\) the
    integrand is \(O(\ln x/x^2)\), while near zero it behaves like
    \((\ln x)/y\); both are integrable. For \(y=0\) it is not integrable.
    If \(y<0\), the denominator vanishes at \(x=\sqrt{-y}\). Unless \(y=-1\),
    the numerator does not vanish there and the singularity is of order
    \(1/(x-\sqrt{-y})\), so the integral diverges. If \(y=-1\), both numerator
    and denominator vanish linearly at \(x=1\), and their quotient has a
    removable singularity. Therefore
    \[
        D(F)=(0,+\infty)\cup\{-1\}.
    \]

    The point \(-1\) is isolated in this domain. Fix \(y_0>0\) and let
    \(y_n\to y_0\). Then
    \[
        f_n(x)=\frac{\ln x}{x^2+y_n}
        \longrightarrow
        \frac{\ln x}{x^2+y_0}
    \]
    pointwise. Eventually \(y_n\geq y_0/2\), so
    \[
        |f_n(x)|
        \leq
        \frac{|\ln x|}{x^2+y_0/2},
    \]
    and the function on the right belongs to
    \(\mathcal L^1((0,+\infty))\). Dominated convergence shows that
    \(F(y_n)\to F(y_0)\). Thus \(F\) is continuous on its domain.
\end{snippetsolution}

Conditionally convergent improper Riemann integrals do not fit this
identification. For example,
\[
    \frac{\sin x}{x}\notin\mathcal L^1((0,+\infty)),
\]
although its improper Riemann integral converges. One sometimes calls this an
improper Lebesgue integral, but it is not a Lebesgue integral in the
\(\mathcal L^1\) sense.

\end{document}
